\section{Methodology}\label{methodology}

\subsection{Problem Formulation}

Let a VQA example be represented as follows:
\[(x,q,y),\]

where $x$ denotes the full image, $q$ the natural-langauge question, and $y=(y_1,\ldots,y_T)$ the 
ground-truth answer sequence. Given a pretrained vision-language model $p_\theta$, the objective
is to adapt the model such that it predicts the correct answer using only the full image and question
at inference.

The two post-training strategies use the same training examples and parameter-efficient
LoRA adaptation \parencite{hu2021lora}. The two strategies primarily differ in their
training objectives and sources of supervision. SFT learns directly from
the ground-truth answer, whereas OPSD learns from the token-level distribution of a privileged teacher
conditioned on question-relevant visual evidence.

\subsection{Privileged Visual Context}

For OPSD, each training example also contains privileged visual input, $x^{p}$. The
GQA dataset provides scene-graph objects and question-grounding information that allows one
to associate questions with the relevant objects in the image. The bounding boxes related to their
relevant objects were combined into a singular region and then expanded slightly to preserve
some surrounding visual context. The outcome is the crops, $x^{p}$. However, it is important to note,
that the crop is based solely on the question and its grounded objects. The ground-truth answer
was not used to select the crop. As a result, the privileged view does not directly reveal the
target label but only provides additional localisation information. The crop therefore does not 
introduce new visual information but instead emphasises question-relevant regions while reducing 
potentially distracting surrounding content.

The student receives the standard inference-time context
\[c_s = (x,q),\]
while the teacher receives
\[c_t = (x^{p},q).\]

In terms of question curation, questions that would result in a local crop removing useful 
information required to answer the question are excluded during dataset construction.

\subsection{Supervised Fine-Tuning (SFT)}

SFT directly optimises the likelihood of the ground-truth answer sequence. For an answer 
$y=y_1,\ldots,y_T$, the autoregressive SFT objective is
\[
-\sum_{t=1}^{T}
\log
p_\theta
\left(
y_t
\mid
y_{<t},x,q
\right).
\]

The model is therefore trained under teacher forcing, where each token is predicted while 
conditioning on the preceding ground-truth tokens $y_{<t}$. At inference time, however, the 
model conditions on its own previously generated tokens. SFT consequently provides strong direct
supervision but trains the model on trajectories that may differ from those encountered during
autoregressive generation.

\subsection{On-Policy Self-Distillation}

OPSD instead trains the student on trajectories generated by its own policy. Given the full-image
context $c_s$, the student first generates an answer trajectory
\[
(\hat{y}_1,\ldots,\hat{y}_K),
\qquad
\hat{y}
\sim
p_\theta(\cdot\mid x,q).
\]

The generated token sequence is then fixed, and both the student and teacher score the 
\emph{same} trajectory. At token position $t$, the student distribution, $p_t$,  is
\[p_t = 
p_\theta
\left(
\cdot
\mid
\hat{y}_{<t},x,q
\right),
\]

while the privileged teacher distribution, $q_t$, is
\[q_t = 
p_{\bar{\theta}}
\left(
\cdot
\mid
\hat{y}_{<t},x^{p},q
\right).
\]

Both distributions assign probabilities over the same vocabulary, conditioned on the 
student's generated tokens, but with different visual inputs. 
Here, $\theta$ denotes the trainable student parameters, while $\bar{\theta}$ denotes the 
frozen teacher parameters. Both originate from the same pretrained VLM, however, only the student
LoRA parameters are updated during OPSD training. Gradients are only propagated through the
student branch \parencite{zhao2601self}. Unlike the EMA-based teacher employed in
Vision-OPD~\cite{yuan2026vision}, the teacher parameters remain fixed throughout training.


The student and teacher distributions are then aligned using the Jensen-Shannon divergence (JSD). For each
generated token position, an intermediate distribution is defined as
\[m_t = 
\frac{1}{2}
\left(
p_t+q_t
\right).
\]

The token-level divergence is then
\[\operatorname{JSD}(p_t\|q_t) = 
\frac{1}{2}
D_{\mathrm{KL}}(p_t|m_t)
+
\frac{1}{2}
D_{\mathrm{KL}}(q_t|m_t).
\]

Here, $D_{\mathrm{KL}}$ denotes the Kullback-Leiber (KL) divergence, which measures the 
discrepancy between two probability distributions:
\[
D_{\mathrm{KL}}(p\|q)
=\sum_{v \in \mathcal{V}}
p(v)\log\frac{p(v)}{q(v)},
\]

where $\mathcal{V}$ represents the token vocabulary. KL divergence is asymmetric and potentially 
unbounded. In contrast, JSD provides a symmetric, bound measure of distributional disagreement, ranging 
from $0$ to $\ln 2$. JSD was therefore selected to provide a balanced objective for aligning 
student and teacher distributions, consistent with Vision-OPD \parencite{yuan2026vision}. Minimising 
JSD encourages the student to approximate the privileged teacher's predictions, although this does not
guarantee their correctness.

The OPSD training objective is
the mean divergence across the student-generated trajectory:
\[\mathcal{L}_{\mathrm{OPSD}} = 
\frac{1}{K}
\sum_{t=1}^{K}
\operatorname{JSD}(p_t|q_t).
\]

Unlike SFT, this loss is evaluated along the student's own generated trajectory rather than a 
fixed reference sequence. The teacher therefore provides dense token-level guidance specifically on the 
states encountered by the student during generation.

It is also important to note that no supervised cross-entropy term is added to the primary OPSD condition. 
This isolates the contribution of privileged-context self-distillation and allows a direct comparison
against conventional SFT. Although the privileged crop is used during OPSD training, it is discarded
during evaluation. All final predictions are generated using only the full image and question, ensuring
SFT and OPSD operate under identical inference conditions.